% cap03/3.2_nivel.tex
\section{Nivel de Investigaci\'on}

La investigaci\'on alcanza un \textbf{nivel descriptivo}. Su prop\'osito es
caracterizar de manera rigurosa y objetiva las propiedades y el comportamiento de
la plataforma \textit{Democra} ---su funcionamiento, su aislamiento entre
inquilinos, su rendimiento y su usabilidad--- sin manipular experimentalmente las
variables ni inducir relaciones de causalidad entre ellas.

\subsection{Alcance de la Descripci\'on}

Describir, en el contexto de la ingenier\'ia de software, significa observar y
registrar de forma sistem\'atica atributos verificables del artefacto: los
c\'odigos de estado y la latencia de la API, la cobertura de c\'odigo alcanzada por
las suites de prueba, la efectividad de las pol\'iticas de seguridad a nivel de
fila y la percepci\'on de usabilidad de los operadores. La descripci\'on se apoya
en evidencia obtenida mediante instrumentos objetivos (pruebas automatizadas,
analizadores de cobertura y la escala estandarizada de usabilidad), evitando
apreciaciones subjetivas no fundamentadas.

De este modo, el nivel descriptivo permite documentar con precisi\'on qu\'e hace el
sistema y con qu\'e grado de calidad lo hace, dejando constancia de sus atributos
para su posterior an\'alisis y discusi\'on.
